% Diagrams for Lesson 9: electromagnetic induction.
\newcommand{\faradayHistoryContext}{%
  \historyfigure[0.93\linewidth]{faraday-ring-built-in-v1.png}%
  {左：Faraday 肖像；右：铁环上的双线圈、检流计与电池构成早期电磁感应实验}}

\newcommand{\movingMagnetExperimentFigure}{%
\begin{center}
\resizebox{0.98\linewidth}{!}{%
\begin{tikzpicture}[x=1cm,y=1cm,>=Latex]
  \node[inner sep=0] at (0,0)
    {\includegraphics[width=14cm]{moving-magnet-experiment-built-in-v1.png}};
  \node[anchor=east,align=center] at (-7.35,2.35)
    {(a) 靠近\\$\vect v_M\leftarrow$};
  \node[anchor=east,align=center] at (-7.35,0)
    {(b) 静止\\$\vect v_M=\vect0$};
  \node[anchor=east,align=center] at (-7.35,-2.35)
    {(c) 远离\\$\vect v_M\rightarrow$};
\end{tikzpicture}}%
\par\smallskip{\small\color{TextGray}图 9.2：磁铁靠近时检流计向一侧偏转；停住时回零；远离时反向偏转。}
\end{center}}

\newcommand{\fluxOrientationDiagram}{%
\begin{center}
\begin{tikzpicture}[scale=0.88,>=Latex]
  \draw[curve] (0,0) ellipse (2.7 and 1.25);
  \draw[->,DeepBlue,very thick] (2.55,0.35) arc[start angle=16,end angle=105,x radius=2.7,y radius=1.25];
  \node[DeepBlue] at (2.5,2.0) {回路正向};
  \draw[vector] (0,0)--(0,2.15) node[above] {$\vect n$};
  \foreach \x in {-1.7,-0.8,0.15,1.1,1.9}{
    \draw[component] (\x,-1.65)--(\x,1.65);
  }
  \node[DeepBlue] at (2.35,0.0) {$\vect B$};
  \node[align=left,anchor=west] at (3.7,0.35)
    {$\displaystyle\Phi_B=\int_S\vect B\cdot\dd\vect S$\\[5pt]
     $\displaystyle\mathcal E=-\frac{\dd\Phi_B}{\dd t}$\\[5pt]
     负号表示感应效应阻碍磁通的改变};
\end{tikzpicture}
\end{center}}

\newcommand{\motionalRodDiagram}{%
\begin{center}
\begin{tikzpicture}[x=1cm,y=1cm,>=Latex]
  \draw[Ink,very thick] (-3,1.5)--(3,1.5);
  \draw[Ink,very thick] (-3,-1.5)--(3,-1.5);
  \draw[Ink,very thick] (-3,-1.5)--(-3,-0.55)
    decorate[decoration={zigzag,segment length=4pt,amplitude=2pt}] {-- (-3,0.55)}--(-3,1.5);
  \node[anchor=west] at (-2.78,0) {$R$};
  \draw[DeepBlue,very thick] (0,-1.65)--(0,1.65);
  \draw[vector] (0.2,0.35)--(1.75,0.35) node[right] {$\vect v$};
  \draw[vector] (-0.18,-0.25)--(-1.45,-0.25) node[left] {$\vect F_B$};
  \draw[vector] (0.18,-0.85)--(1.45,-0.85) node[right] {$\vect F_{\rm ext}$};
  \draw[<->] (-3.55,-1.5)--node[left] {$\ell$}(-3.55,1.5);
  \foreach \x in {-2.1,-1.3,0.7,1.5}{
    \foreach \y in {-1.2,1.15}{\node[TextGray] at (\x,\y) {$\times$};}
  }
  \node[DeepBlue] at (2.55,1.0) {$\vect B$};
  \draw[vector] (0,-0.82)--(0,0.82) node[left] {$I$};
  \draw[->,DeepBlue,thick] (-0.55,1.5)--(-1.65,1.5) node[midway,above] {$I$};
  \draw[->,DeepBlue,thick] (-1.65,-1.5)--(-0.55,-1.5);
  \node[align=left,anchor=west] at (3.4,0.3)
    {$\displaystyle\mathcal E_{\rm m}=\int(\vect v\times\vect B)\cdot\dd\vect\ell$\\[5pt]
     $\displaystyle\mathcal E_{\rm m}=B\ell v$ （三者两两垂直）\\[5pt]
     $\displaystyle I=\frac{B\ell v}{R},\quad F_B=BI\ell$\\[4pt]
     $q>0:\ q(\vect v\times\vect B)$ 沿棒向上};
\end{tikzpicture}
\end{center}}

\newcommand{\movingLoopDiagram}{%
\begin{center}
\begin{tikzpicture}[scale=0.9,>=Latex]
  \draw[curve,postaction={decorate},decoration={markings,mark=at position .38 with {\arrow{Latex}}}]
    plot[smooth cycle] coordinates{(-2,-0.7)(-1,1.2)(0.8,1.0)(2.0,0)(1.1,-1.1)(-0.8,-1.3)};
  \node[DeepBlue,anchor=south west] at (0.45,1.12) {$C$ 正向};
  \draw[vector] (-1.2,-0.5)--(-0.25,0.15) node[above] {$\vect v$};
  \foreach \p in {(-1.5,0.3),(-0.5,0.7),(0.45,0.35),(1.05,-0.45),(-0.2,-0.75)}
    {\node[TextGray] at \p {$\times$};}
  \node[DeepBlue] at (2.5,-0.4) {$\vect B$};
  \node[align=left,anchor=west] at (3.4,0.2)
   {$\displaystyle\mathcal E=\oint_C(\vect E+\vect v\times\vect B)\cdot\dd\vect\ell$\\[5pt]
    $\displaystyle\mathcal E=-\frac{\dd}{\dd t}\int_{S(t)}\vect B\cdot\dd\vect S$\\[5pt]
    图示特例：匀强 $\vect B$ 中刚性平移，$\mathcal E=0$};
\end{tikzpicture}
\end{center}}

\newcommand{\vortexFieldDiagram}{%
\begin{center}
\begin{tikzpicture}[scale=0.88,>=Latex]
  \draw[Ink,thick] (0,0) circle (2.0);
  \foreach \r in {0.65,1.25,2.65,3.35}{
    \draw[DeepBlue,thick,->] (\r,0) arc[start angle=0,end angle=180,radius=\r];
    \draw[DeepBlue,thick] (-\r,0) arc[start angle=180,end angle=360,radius=\r];
  }
  \foreach \x in {-1.1,0,1.1}{\foreach \y in {-1.0,0,1.0}{
    \pgfmathparse{\x*\x+\y*\y<3.5?1:0}\ifnum\pgfmathresult=1\node[TextGray] at (\x,\y) {$\times$};\fi}}
  \draw[<->] (0,-3.85)--node[below] {$R$}(2.0,-3.85);
  \node[align=left,anchor=west] at (4.1,0.25)
    {$\displaystyle E_\varphi(r)=-\frac r2\dot B,\quad r\le R$\\[6pt]
     $\displaystyle E_\varphi(r)=-\frac{R^2}{2r}\dot B,\quad r\ge R$\\[4pt]
     图示 $\dot B>0$ 向纸内；$\varphi$ 正向为顺时针\\感生电场沿逆时针};
\end{tikzpicture}
\end{center}}

\newcommand{\selfInductanceDiagram}{%
\begin{center}
\begin{tikzpicture}[scale=0.86,>=Latex]
  \foreach \x in {-2.4,-2.0,...,2.4}{\draw[Ink,thick] (\x,0) ellipse (0.22 and 1.45);}
  \draw[DeepBlue,->,very thick] (-3.0,0)--(3.0,0) node[right] {$\vect B$};
  \draw[Ink,very thick,->] ({2.4+.22*cos(310)},{1.45*sin(310)}) arc(310:380:.22 and 1.45);
  \node[right] at (2.7,-.6) {$I$};
  \draw[guide,dashed] (-2.6,-1.75) rectangle (2.6,1.75);
  \node[align=left,anchor=west] at (4.0,0.25)
    {$\displaystyle\Lambda=N\Phi=LI$\\[5pt]
     $\displaystyle\mathcal E_L=-\frac{\dd(LI)}{\dd t}$\\[5pt]
     $\displaystyle W_L=\frac12LI^2$};
\end{tikzpicture}
\end{center}}

\newcommand{\rlInductorContext}{%
\begin{center}
\resizebox{0.96\linewidth}{!}{%
\begin{tikzpicture}[x=1cm,y=1cm,>=Latex,line cap=round,line join=round]
  \node[inner sep=0] at (-3.75,0)
    {\includegraphics[width=5.4cm]{rl-inductor-built-in-v1.png}};
  \begin{scope}[shift={(0.1,0)}]
    \draw[Ink,thick] (0,1.55)--(1.35,1.55);
    \draw[Ink,thick,decorate,decoration={zigzag,segment length=5pt,amplitude=2.5pt}]
      (1.35,1.55)--(3.15,1.55);
    \draw[Ink,thick] (3.15,1.55)--(5.1,1.55)--(5.1,0.72);
    \foreach \yy in {0.54,0.18,-0.18,-0.54}{
      \draw[Ink,thick] (5.1,{\yy+0.18}) arc[start angle=90,end angle=-90,x radius=0.25,y radius=0.18];
    }
    \draw[Ink,thick] (5.1,-0.72)--(5.1,-1.55)--(3.20,-1.55);
    \fill[white,draw=Ink,thick] (3.20,-1.55) circle (2.2pt);
    \fill[white,draw=Ink,thick] (2.25,-1.55) circle (2.2pt);
    \draw[Ink,thick] (2.25,-1.55)--(3.03,-1.10);
    \draw[Ink,thick] (2.25,-1.55)--(0,-1.55)--(0,-0.30);
    \draw[Ink,thick] (0,0.30)--(0,1.55);
    \draw[Ink,very thick] (-0.38,0.30)--(0.38,0.30);
    \draw[Ink,very thick] (-0.23,-0.30)--(0.23,-0.30);
    \node[left] at (-0.42,0.30) {$+$};
    \node[left] at (-0.30,-0.30) {$-$};
    \node[left] at (-0.48,0) {$\mathcal E$};
    \node[above] at (2.25,1.72) {$R$};
    \node[right] at (5.38,0) {$L$};
    \node[above left=3pt] at (5.12,0.72) {$a$};
    \node[below left=3pt] at (5.12,-0.72) {$b$};
    \node[below] at (2.67,-1.78) {$S$};
    \draw[->,DeepBlue,very thick] (0.55,1.90)--(1.35,1.90) node[midway,above] {$I$};
  \end{scope}
\end{tikzpicture}}%
\par\smallskip{\small\color{TextGray}图 9.9：左为环形电感实物结构；右为串联 $RL$ 回路。闭合 $S$ 后，电流沿箭头方向逐渐建立。}
\end{center}}

\newcommand{\mutualInductanceDiagram}{%
\begin{center}
\begin{tikzpicture}[scale=0.9,>=Latex]
  \draw[curve] (-2.0,-0.25) ellipse (2.3 and 1.1);
  \draw[curve] (2.0,0.45) ellipse (2.0 and 0.95);
  \draw[->,DeepBlue,very thick] (-0.1,0.15) arc[start angle=10,end angle=70,x radius=2.3,y radius=1.1];
  \draw[->,DeepBlue,very thick] (3.75,0.7) arc[start angle=15,end angle=70,x radius=2.0,y radius=0.95];
  \node at (-2.0,-0.25) {$1$}; \node at (2.0,0.45) {$2$};
  \node[DeepBlue] at (-0.1,1.25) {$I_1$}; \node[DeepBlue] at (3.75,1.4) {$I_2$};
  \foreach \x in {-0.7,0.1,0.9}{\draw[component] (\x,-1.5) to[bend left=18] (\x+0.9,1.8);}
  \node[align=left,anchor=west] at (5.0,0.45)
   {$\Phi_2=M_{12}I_1,\quad\Phi_1=M_{21}I_2$\\[5pt]
    $\boxed{M_{12}=M_{21}}$\\[5pt]
    $W=\tfrac12L_1I_1^2+MI_1I_2+\tfrac12L_2I_2^2$};
\end{tikzpicture}
\end{center}}

\newcommand{\concentricLoopsDiagram}{%
\begin{center}\begin{tikzpicture}[scale=0.78,>=Latex]
  \draw[Ink,very thick] (0,0) circle (2.0); \draw[DeepBlue,very thick] (0,0) circle (1.1);
  \draw[->,Ink,thick] (1.85,0.75) arc[start angle=22,end angle=-45,radius=2.0];
  \draw[->,DeepBlue,thick] (0.95,0.55) arc[start angle=30,end angle=120,radius=1.1];
  \node[Ink,anchor=west] at (2.65,0.4) {原电流}; \node[DeepBlue] at (0,1.55) {感应电流};
\end{tikzpicture}\end{center}}

\newcommand{\magnetThreeLoopsDiagram}{%
\begin{center}\begin{tikzpicture}[scale=0.8,>=Latex]
  \draw[DeepBlue,dashed,very thick] (0,1.25) arc(90:270:.33 and 1.25);
  \draw[fill=SoftGray] (-1.8,-0.65) rectangle (1.8,0.65);
  \node at (-1.25,0) {$S$}; \node at (1.25,0) {$N$};
  \foreach \x in {-3.4,3.4}{\draw[DeepBlue,very thick] (\x,0) ellipse (0.33 and 1.25);}
  \draw[DeepBlue,very thick] (0,-1.25) arc(270:450:.33 and 1.25);
  \draw[vector] (-1.7,1.1)--(-2.7,1.1) node[midway,above=4pt] {$\vect v_M$};
  \draw[vector] (-3.4,-1.55)--(-4.45,-1.55) node[midway,below] {$\vect a_1$};
  \node[anchor=north] at (0,-1.43) {$t=0:\ \vect a_2=0$};
  \draw[vector,dashed] (0,-2.23)--(-1.15,-2.23)
    node[midway,below] {左移后 $\vect a_2$};
  \draw[vector] (3.4,-1.55)--(2.35,-1.55) node[midway,below] {$\vect a_3$};
  \node[TextGray,align=center] at (-3.4,2.1) {极外近轴\\收缩};
  \node[TextGray,align=center] at (0,3.5) {初始无趋势\\近中部左移后收缩};
  \node[TextGray,align=center] at (3.4,2.1) {极外近轴\\扩张};
\end{tikzpicture}\end{center}}

\newcommand{\rodResistorDiagram}{%
\begin{center}\begin{tikzpicture}[scale=0.78,>=Latex]
  \draw[Ink,very thick] (-3,1.4)--(3,1.4); \draw[Ink,very thick] (-3,-1.4)--(3,-1.4);
  \draw[Ink,very thick] (-2.7,-1.4)--(-2.7,-0.45) decorate[decoration={zigzag,segment length=4pt,amplitude=2pt}] {-- (-2.7,0.45)}--(-2.7,1.4);
  \node[left=8pt] at (-2.7,0) {$R$};
  \draw[DeepBlue,very thick] (0.8,-1.55)--(0.8,1.55); \node[above] at (0.8,1.55) {$m$};
  \draw[vector] (1.0,0)--(2.4,0) node[right] {$v$};
  \draw[<->] (-3.7,-1.4)--node[left] {$L$}(-3.7,1.4);
  \foreach \x in {-1.7,-0.7,0.3,1.4,2.3}{\foreach \y in {-0.75,0.75}{\node[TextGray] at (\x,\y) {$\times$};}}
  \node[align=left,anchor=west] at (3.7,0.2)
    {$\displaystyle I=\frac{BLv}{R}$\\[4pt]
     $\displaystyle m\dot v=-\frac{B^2L^2}{R}v$};
\end{tikzpicture}\end{center}}

\newcommand{\twoRodDiagram}{%
\begin{center}\begin{tikzpicture}[scale=0.78,>=Latex]
  \draw[Ink,very thick] (-3,1.4)--(3,1.4); \draw[Ink,very thick] (-3,-1.4)--(3,-1.4);
  \draw[DeepBlue,very thick] (-1.5,-1.55)--(-1.5,1.55); \draw[DeepBlue,very thick] (1.5,-1.55)--(1.5,1.55);
  \node[above] at (-1.5,1.62) {$m,R$}; \node[above] at (1.5,1.62) {$m,R$};
  \draw[vector] (1.7,0)--(2.9,0) node[right] {$v_0$};
  \foreach \x in {-2.4,-0.5,0.5,2.4}{\foreach \y in {-0.7,0.7}{\node[TextGray] at (\x,\y) {$\times$};}}
  \draw[<->] (-3.25,-1.4)--node[left] {$L$}(-3.25,1.4);
  \node[align=left,anchor=west] at (3.7,0.1)
   {$v_{\rm cm}=v_0/2$\\[4pt]$u=v_2-v_1=v_0\e^{-t/\tau}$};
\end{tikzpicture}\end{center}}

\newcommand{\unequalRodDiagram}{%
\begin{center}\begin{tikzpicture}[x=0.85cm,y=0.85cm,>=Latex]
  \draw[Ink,very thick] (-5,1.0)--(-1,1.0); \draw[Ink,very thick] (-5,-1.0)--(-1,-1.0);
  \draw[DeepBlue,very thick] (-4.0,-1.15)--(-4.0,1.15) node[above] {$m_1$};
  \draw[Ink,dotted,very thick] (-1,0)--(1,0);
  \draw[Ink,very thick] (1,1.8)--(5,1.8); \draw[Ink,very thick] (1,-1.8)--(5,-1.8);
  \draw[DeepBlue,very thick] (4.0,-1.95)--(4.0,1.95) node[above] {$m_2$};
  \draw[vector] (4.2,0)--(5.5,0) node[right] {$v_0$};
  \node at (-2.2,0) {$B_1,L_1$}; \node at (2.2,0) {$B_2,L_2$};
  \foreach \x in {-4.7,-3.2,-1.7,1.5,3.0,4.6}{\node[TextGray] at (\x,0.55) {$\times$};}
\end{tikzpicture}\end{center}}

\newcommand{\rodCapacitorDiagram}{%
\begin{center}\begin{tikzpicture}[scale=0.78,>=Latex]
  \draw[Ink,very thick] (-3,1.4)--(3,1.4); \draw[Ink,very thick] (-3,-1.4)--(3,-1.4);
  \draw[Ink,very thick] (-2.7,-1.4)--(-2.7,-0.25); \draw[Ink,very thick] (-3.05,-0.25)--(-2.35,-0.25);
  \draw[Ink,very thick] (-3.05,0.25)--(-2.35,0.25); \draw[Ink,very thick] (-2.7,0.25)--(-2.7,1.4);
  \node[left] at (-2.7,0) {$C$};
  \draw[DeepBlue,very thick] (0.8,-1.55)--(0.8,1.55) node[above] {$m$};
  \draw[vector] (1.0,0)--(2.4,0) node[right] {$v_0$};
  \foreach \x in {-1.7,-0.7,0.2,1.5,2.4}{\foreach \y in {-0.7,0.7}{\node[TextGray] at (\x,\y) {$\times$};}}
  \node[align=left,anchor=west] at (3.5,0.2)
    {$q_\infty=CBLv_\infty$\\[4pt]$m(v-v_0)=-BLq$};
\end{tikzpicture}\end{center}}

\newcommand{\rodInductorDiagram}{%
\begin{center}\begin{tikzpicture}[scale=0.78,>=Latex]
  \draw[Ink,very thick] (-3,1.4)--(3,1.4); \draw[Ink,very thick] (-3,-1.4)--(3,-1.4);
  \draw[Ink,very thick] (-2.7,-1.4)--(-2.7,-0.6);
  \draw[Ink,thick] (-2.7,-0.6) decorate[decoration={coil,aspect=0.65,segment length=5pt,amplitude=5pt}] {-- (-2.7,0.6)};
  \draw[Ink,very thick] (-2.7,0.6)--(-2.7,1.4); \node[left] at (-2.9,0) {$\mathcal L$};
  \draw[DeepBlue,very thick] (0.8,-1.55)--(0.8,1.55) node[above] {$m$};
  \draw[vector] (1.0,0)--(2.4,0) node[right] {$v_0$};
  \foreach \x in {-1.7,-0.7,0.2,1.5,2.4}{\foreach \y in {-0.7,0.7}{\node[TextGray] at (\x,\y) {$\times$};}}
  \node[align=left,anchor=west] at (3.5,0.2)
    {$\mathcal L\dot I=B\ell v$\\[4pt]$m\dot v=-B\ell I$};
\end{tikzpicture}\end{center}}

\newcommand{\rodSourceDiagram}{%
\begin{center}\begin{tikzpicture}[scale=0.78,>=Latex]
  \draw[Ink,very thick] (-3,1.4)--(3,1.4); \draw[Ink,very thick] (-3,-1.4)--(3,-1.4);
  \draw[Ink,very thick] (-2.7,-1.4)--(-2.7,-0.55);
  \draw[Ink,thick] (-2.7,-0.55) decorate[decoration={zigzag,segment length=4pt,amplitude=2pt}] {-- (-2.7,0.1)};
  \draw[Ink,very thick] (-3.05,0.25)--(-2.35,0.25); \draw[Ink,thick] (-2.9,0.55)--(-2.5,0.55);
  \draw[Ink,very thick] (-2.7,0.55)--(-2.7,1.4); \node[left] at (-2.85,-0.25) {$R$}; \node[left] at (-2.9,0.5) {$\mathcal E$};
  \draw[DeepBlue,very thick] (0.8,-1.55)--(0.8,1.55) node[above] {$m$};
  \draw[vector] (1.0,0)--(2.4,0) node[right] {$v$};
  \foreach \x in {-1.7,-0.7,0.2,1.5,2.4}{\foreach \y in {-0.7,0.7}{\node[TextGray] at (\x,\y) {$\times$};}}
\end{tikzpicture}\end{center}}

\newcommand{\equivalentCircuitDiagram}{%
\begin{center}\begin{tikzpicture}[scale=0.82,>=Latex]
  \draw[Ink,very thick] (-4,0)--(-2.7,0);
  \draw[Ink,thick] (-2.7,0) decorate[decoration={zigzag,segment length=5pt,amplitude=2.5pt}] {-- (-0.9,0)};
  \node[above] at (-1.8,0.25) {$R$};
  \draw[Ink,very thick] (-0.9,0)--(0.0,0);
  \draw[Ink,very thick] (0,-0.75)--(0,0.75); \draw[Ink,very thick] (0.45,-0.75)--(0.45,0.75);
  \node[above] at (0.22,0.9) {$C_m=m/(B^2L^2)$};
  \draw[Ink,very thick] (0.45,0)--(2.0,0);
  \node[align=left,anchor=west] at (3.0,0.2)
   {$U=RI+U_m$\\[4pt]$I=C_m\dot U_m$\\[4pt]$\tfrac12mv^2=\tfrac12C_mU_m^2$};
\end{tikzpicture}\end{center}}

\newcommand{\chargedRingDiagram}{%
\begin{center}\begin{tikzpicture}[scale=0.82,>=Latex]
  \draw[Ink,very thick] (0,0) circle (2.4);
  \fill[DeepBlue] (2.4,0) circle (3pt);
  \node[anchor=south west] at (2.5,0.12) {$m,+q$};
  \draw[->,DeepBlue,very thick] (2.0,1.15) arc[start angle=30,end angle=105,radius=2.4];
  \node[DeepBlue] at (2.8,1.5) {$\dot\theta$};
  \foreach \x in {-1.1,0,1.1}{\foreach \y in {-1.1,0,1.1}{\node[TextGray] at (\x,\y) {$\times$};}}
  \node at (0,0) {$B(t)$};
  \draw[vector] (2.32,-0.32)--(1.05,-0.32)
    node[midway,below] {$a_r=-R\dot\theta^2$};
  \draw[axis] (2.48,0.32)--(3.75,0.32)
    node[right,align=left] {径向正向\\$N_r,\ qR\dot\theta B$ 均取代数值};
\end{tikzpicture}\end{center}}

\newcommand{\chordParticleDiagram}{%
\begin{center}\begin{tikzpicture}[scale=0.8,>=Latex]
  \draw[Ink,thick] (0,0) circle (2.5);
  \draw[Ink,very thick] ({-sqrt(2.5*2.5-1.25*1.25)},-1.25)--({sqrt(2.5*2.5-1.25*1.25)},-1.25);
  \fill[DeepBlue] (0,-1.25) circle (3pt) node[below] {$m,+q$};
  \draw[<->] (0,0)--node[right] {$h$}(0,-1.25);
  \draw[axis] (-2.2,-1.75)--(2.4,-1.75) node[right] {$x$};
  \node[TextGray] at (-1.4,1.0) {$\times$}; \node[TextGray] at (0,1.0) {$\times$}; \node[TextGray] at (1.4,1.0) {$\times$};
  \node at (0,1.65) {$B_0\sin\omega t$};
  \node[align=left,anchor=west] at (3.5,0)
    {$\displaystyle m\ddot x=\frac{qh}{2}\dot B$\\[5pt]
     $\displaystyle x=\frac{qhB_0}{2m\omega}(1-\cos\omega t)$};
\end{tikzpicture}\end{center}}

\newcommand{\chordEmfDiagram}{%
\begin{center}\begin{tikzpicture}[scale=0.82,>=Latex]
  \draw[Ink,thick] (0,0) circle (2.5);
  \draw[Ink,very thick] (-3.2,-1.2)--(4.3,-1.2);
  \coordinate (A) at ({-sqrt(2.5*2.5-1.2*1.2)},-1.2); \coordinate (C) at ({sqrt(2.5*2.5-1.2*1.2)},-1.2); \coordinate (D) at (4.1,-1.2);
  \fill[DeepBlue] (A) circle (2.5pt) node[below] {$A$}; \fill[DeepBlue] (C) circle (2.5pt) node[below] {$C$}; \fill[Ink] (D) circle (2.5pt) node[below] {$D$};
  \draw[<->] ({sqrt(2.5*2.5-1.2*1.2)},-2)--(4.1,-2) node[midway,below] {$\ell$};
  \draw[<->] (0,0)--node[right] {$h$}(0,-1.2);
  \foreach \r in {0.8,1.6,3.2}{\draw[DeepBlue,thick,->] (\r,0) arc[start angle=0,end angle=105,radius=\r];}
  \node[TextGray] at (-1.1,.2) {$\dot B=k$};
\end{tikzpicture}\end{center}}

\newcommand{\circularVoltmetersDiagram}{%
\begin{center}\begin{tikzpicture}[scale=0.76,>=Latex]
  \draw[Ink,very thick] (0,0) circle (2.4);
  \coordinate (A) at (150:2.4); \coordinate (C) at (30:2.4);
  \fill[DeepBlue] (A) circle (2.5pt) node[left] {$A$}; \fill[DeepBlue] (C) circle (2.5pt) node[right] {$C$};
  \draw[Ink,thick] (A) to[out=115,in=65,looseness=1.05]
    node[midway,circle,draw,fill=white,inner sep=2pt] {$V_1$} (C);
  \draw[Ink,thick] (A) .. controls (-3.2,-2.8) and (3.2,-2.8) ..
    node[pos=0.5,circle,draw,fill=white,inner sep=2pt] {$V_2$} (C);
  \draw[Ink,thick] (A)--(-0.35,0.2);
  \draw[Ink,thick] (0.35,0.2)--(C);
  \node[circle,draw,fill=white,inner sep=2pt] at (0,0.2) {$V_3$};
  \foreach \x in {-1,0,1}{\node[TextGray] at (\x,-0.65) {$\times$};}
  \node at (0,-1.15) {$\dot B=k$};
\end{tikzpicture}\end{center}}

\newcommand{\squareVoltmetersDiagram}{%
\begin{center}\begin{tikzpicture}[scale=0.76,>=Latex]
  \draw[Ink,very thick] (-2,-2) rectangle (2,2);
  \draw[Ink,thick] (-2,2.55)--(2,2.55); \node[circle,draw,fill=white,inner sep=2pt] at (0,2.55) {$V_1$};
  \draw[Ink,thick] (2.55,-2)--(2.55,2); \node[circle,draw,fill=white,inner sep=2pt] at (2.55,0) {$V_2$};
  \draw[Ink,thick] (-2,-2.55)--(2,-2.55); \node[circle,draw,fill=white,inner sep=2pt] at (0,-2.55) {$V_3$};
  \draw[Ink,thick] (-2.55,-2)--(-2.55,2); \node[circle,draw,fill=white,inner sep=2pt] at (-2.55,0) {$V_4$};
  \node[left] at (-2,0.75) {$R_1$}; \node[right] at (2,0.75) {$R_2$};
  \node[TextGray] at (0,0) {$\times\ B(t)$}; \draw[<->] (-2,-2.85)--node[below] {$a$}(2,-2.85);
\end{tikzpicture}\end{center}}

\newcommand{\gridCurrentDiagram}{%
\begin{center}\begin{tikzpicture}[scale=0.88,>=Latex]
  \draw[Ink,very thick] (-2,-2) grid[step=2] (2,2);
  \draw[->,DeepBlue,very thick] (-2,2)--(0,2);
  \draw[->,DeepBlue,very thick] (0,2)--(2,2);
  \draw[->,DeepBlue,very thick] (2,2)--(2,0);
  \draw[->,DeepBlue,very thick] (2,0)--(2,-2);
  \draw[->,DeepBlue,very thick] (2,-2)--(0,-2);
  \draw[->,DeepBlue,very thick] (0,-2)--(-2,-2);
  \draw[->,DeepBlue,very thick] (-2,-2)--(-2,0);
  \draw[->,DeepBlue,very thick] (-2,0)--(-2,2);
  \node[DeepBlue] at (2.6,0) {$I_0$};
  \node[TextGray] at (0,0.35) {内部支路}; \node[TextGray] at (0,-0.35) {$I=0$};
  \draw[<->] (-2,-2.5)--node[below] {$2a$}(2,-2.5);
\end{tikzpicture}\end{center}}

\newcommand{\solenoidSplitDiagram}{%
\begin{center}\begin{tikzpicture}[scale=0.78,>=Latex]
  \foreach \x in {-4.0,-3.6,...,4.0}{\draw[Ink,thick] (\x,0) ellipse (0.2 and 1.15);}
  \draw[guide,dashed] (0,-1.55)--(0,1.55);
  \draw[<->] (-4.2,-1.65)--node[below] {$l_1$}(0,-1.65);
  \draw[<->] (0,-1.65)--node[below] {$l_2$}(4.2,-1.65);
  \draw[vector] (-3.7,0)--(3.7,0) node[right] {$B=\mu_0nI$};
  \node[align=left,anchor=west] at (5.0,0.15)
    {$L=\mu_0n^2A(l_1+l_2)$\\[4pt]$=L_1+L_2$ （理想长螺线管）};
\end{tikzpicture}\end{center}}

\newcommand{\mutualLoopsProofDiagram}{%
\begin{center}\begin{tikzpicture}[scale=0.82,>=Latex]
  \draw[curve] (-2.6,0) ellipse (1.35 and 0.9); \draw[curve] (2.6,0.4) ellipse (1.25 and 0.8);
  \draw[vector] (-1.5,0.35)--(1.9,0.7) node[midway,above] {$\vect r_{12}$};
  \draw[vector] (-3.2,-0.65)--(-2.2,-0.85) node[below] {$\dd\vect\ell_1$};
  \draw[vector] (2.0,-0.15)--(3.0,-0.25) node[below] {$\dd\vect\ell_2$};
  \node[align=left,anchor=west] at (4.6,0.25)
   {$\displaystyle M_{12}=\frac{\mu_0}{4\pi}\oint_{C_1}\!\oint_{C_2}
   \frac{\dd\vect\ell_1\cdot\dd\vect\ell_2}{r_{12}}$\\[6pt]
    交换 $1,2$ 后积分核不变};
\end{tikzpicture}\end{center}}

\newcommand{\rlcDiagram}[1]{%
\begin{center}\begin{tikzpicture}[scale=0.70,>=Latex]
  \draw[Ink,very thick] (-4,0.9)--(-3.2,0.9);
  \draw[Ink,very thick] (-3.2,0.3)--(-3.2,1.5); \draw[Ink,very thick] (-2.75,0.5)--(-2.75,1.3);
  \node[above] at (-3.0,1.5) {$C,\ U$};
  \draw[Ink,very thick] (-2.75,0.9)--(-1.7,0.9);
  \draw[Ink,thick] (-1.7,0.9) decorate[decoration={coil,aspect=0.65,segment length=5pt,amplitude=5pt}] {-- (0.3,0.9)};
  \node[above] at (-0.7,1.45) {$L$};
  \draw[Ink,thick] (0.3,0.9) decorate[decoration={zigzag,segment length=5pt,amplitude=2.5pt}] {-- (2.2,0.9)};
  \node[above] at (1.25,1.25) {$R$};
  \draw[Ink,very thick] (2.2,0.9)--(3.0,0.9)--(3.0,-0.9);
  \ifnum#1=1
    \draw[Ink,very thick] (3.0,-0.9)--(0.25,-0.9);
    \draw[Ink,very thick] (-0.2,-0.9)--(-4,-0.9)--(-4,0.9);
    \draw[Ink,very thick] (-0.2,-1.4)--(-0.2,-0.4); \draw[Ink,very thick] (0.25,-1.2)--(0.25,-0.6);
    \node[below] at (0,-1.35) {$\mathcal E$};
  \else
    \draw[Ink,very thick] (3.0,-0.9)--(-4,-0.9)--(-4,0.9);
  \fi
  \draw[->,DeepBlue,very thick] (-1.6,0.38)--(0.8,0.38) node[midway,above] {$I(t)$};
\end{tikzpicture}\end{center}}

\newcommand{\travelingFieldTrainDiagram}{%
\begin{center}\begin{tikzpicture}[scale=0.72,>=Latex]
  \draw[Ink,very thick] (-4,1.25)--(4,1.25); \draw[Ink,very thick] (-4,-1.25)--(4,-1.25);
  \draw[DeepBlue,very thick,fill=PaleBlue] (-1.4,-1.25) rectangle (1.4,1.25);
  \node at (0,0) {$m$}; \draw[<->] (-1.4,-1.62)--node[below] {$a$}(1.4,-1.62);
  \draw[<->] (-3.4,-1.25)--node[left] {$L$}(-3.4,1.25);
  \draw[vector] (1.55,0.45)--(3.0,0.45) node[right] {$v$};
  \foreach \x in {-3.2,-2.4,-0.8,0,0.8,2.4,3.2}{\node[TextGray] at (\x,0.75) {$\times$};}
  \draw[->,TextGray,thick] (-3.5,-0.65)--(3.5,-0.65)
    node[right] {$v_\phi$};
  \node[align=left,anchor=west] at (4.6,0.2)
    {$B=B_0\cos(kx-\omega t)$\\[4pt]$v_\phi=\omega/k$};
\end{tikzpicture}\end{center}}
